%=======================================================================
%  Sesión 12 — Virtualización y contenedores en Linux
%  Sistemas Operativos Aplicados a las Telecomunicaciones · UNMSM
%  Compilar con: pdflatex (dos pasadas)
%=======================================================================
\documentclass[11pt,a4paper]{article}

%------------------------- Idioma y codificación -----------------------
\usepackage[T1]{fontenc}
\usepackage[utf8]{inputenc}
\usepackage[spanish,es-noquoting,es-tabla]{babel}

%------------------------- Tipografía Times ----------------------------
\usepackage{newtxtext}      % Times New Roman (texto)
\usepackage{newtxmath}      % Times (matemáticas)
\usepackage[scaled=0.82]{beramono} % monoespaciada legible para comandos

%------------------------- Geometría y espaciado -----------------------
\usepackage[a4paper,top=2.4cm,bottom=2.4cm,left=2.6cm,right=2.6cm]{geometry}
\usepackage{microtype}
\usepackage{setspace}
\setstretch{1.08}
\setlength{\parskip}{0.45em}
\setlength{\parindent}{0pt}
\usepackage{enumitem}
\setlist{itemsep=2pt,topsep=3pt,parsep=0pt}

%------------------------- Color, cajas, gráficos ----------------------
\usepackage{xcolor}
\definecolor{azulprof}{HTML}{14354B}
\definecolor{azulmedio}{HTML}{2E6E8E}
\definecolor{tealx}{HTML}{1A8F86}
\definecolor{vinox}{HTML}{B33A52}
\definecolor{ambarx}{HTML}{C8891B}
\definecolor{grisfondo}{HTML}{F2F5F8}
\definecolor{azulsuave}{HTML}{E6EFF7}
\definecolor{verdesuave}{HTML}{E2F2EE}
\definecolor{ambarsuave}{HTML}{FDF3DF}
\definecolor{rosasuave}{HTML}{FBE8EC}
\definecolor{grislinea}{HTML}{C9D3DC}

\usepackage[most]{tcolorbox}
\usepackage{tikz}
\usetikzlibrary{arrows.meta,positioning,calc,fit,backgrounds,shapes.geometric}
\usepackage{booktabs}
\usepackage{array}
\usepackage{tabularx}
\usepackage{graphicx}
\usepackage{textcomp}

%------------------------- Títulos -------------------------------------
\usepackage{titlesec}
\titleformat{\section}
  {\normalfont\large\bfseries\color{azulprof}}{\thesection.}{0.6em}{}
  [\vspace{2pt}{\color{grislinea}\titlerule[0.9pt]}]
\titleformat{\subsection}
  {\normalfont\normalsize\bfseries\color{azulmedio}}{\thesubsection}{0.6em}{}
\titlespacing*{\section}{0pt}{1.6em}{0.7em}
\titlespacing*{\subsection}{0pt}{1.1em}{0.4em}

%------------------------- Encabezado y pie ----------------------------
\usepackage{fancyhdr}
\pagestyle{fancy}
\fancyhf{}
\renewcommand{\headrulewidth}{0.4pt}
\fancyhead[L]{\footnotesize\color{azulmedio}Sesión 12 · Virtualización y contenedores en Linux}
\fancyfoot[C]{\footnotesize\color{azulmedio}\thepage}

%------------------------- Enlaces -------------------------------------
\usepackage[hidelinks]{hyperref}
\usepackage{url}
\urlstyle{sf}

%------------------------- Cajas reutilizables -------------------------
\newtcolorbox{clave}[1][]{
  enhanced, breakable, colback=azulsuave, colframe=azulmedio,
  boxrule=0pt, leftrule=3pt, arc=2pt, left=8pt, right=8pt, top=6pt, bottom=6pt,
  fontupper=\normalsize, #1}

\newtcolorbox{aviso}[1][]{
  enhanced, breakable, colback=ambarsuave, colframe=ambarx,
  boxrule=0pt, leftrule=3pt, arc=2pt, left=8pt, right=8pt, top=6pt, bottom=6pt, #1}

\newtcolorbox{errorbox}[1][]{
  enhanced, breakable, colback=rosasuave, colframe=vinox,
  boxrule=0pt, leftrule=3pt, arc=2pt, left=8pt, right=8pt, top=6pt, bottom=6pt, #1}

\newtcolorbox{defbox}[1][]{
  enhanced, breakable, colback=verdesuave, colframe=tealx,
  boxrule=0pt, leftrule=3pt, arc=2pt, left=8pt, right=8pt, top=6pt, bottom=6pt, #1}

\tcbuselibrary{listings}
\newtcblisting{codigo}[1][]{
  listing only, colback=azulprof, colframe=azulprof,
  arc=3pt, left=10pt, right=10pt, top=8pt, bottom=8pt, boxrule=0pt,
  listing options={basicstyle=\ttfamily\small\color{white},
    columns=flexible, keepspaces=true, breaklines=true,
    showstringspaces=false, literate={á}{{\'a}}1 {é}{{\'e}}1 {í}{{\'i}}1
      {ó}{{\'o}}1 {ú}{{\'u}}1 {ñ}{{\~n}}1 {µ}{{u}}1},
  #1}

% Etiqueta de bloque
\newcommand{\bloque}[2]{%
  \vspace{1.2em}
  \begin{tcolorbox}[enhanced, colback=azulprof, colframe=azulprof, arc=3pt,
      left=12pt,right=12pt,top=9pt,bottom=9pt]
    {\color{white!75!azulmedio}\footnotesize\bfseries\MakeUppercase{#1}}\\[2pt]
    {\color{white}\large\bfseries #2}
  \end{tcolorbox}
  \vspace{0.2em}}

% Término destacado en línea
\newcommand{\cmd}[1]{\texttt{#1}}

%------------------------- Estilos TikZ --------------------------------
\tikzset{
  caja/.style   = {rectangle, rounded corners=3pt, draw=grislinea, fill=grisfondo,
                   align=center, inner sep=6pt, font=\small, minimum height=9mm},
  cazul/.style  = {caja, fill=azulsuave, draw=azulmedio!40},
  cverde/.style = {caja, fill=verdesuave, draw=tealx!50},
  cambar/.style = {caja, fill=ambarsuave, draw=ambarx!50},
  crosa/.style  = {caja, fill=rosasuave, draw=vinox!40},
  cosc/.style   = {caja, fill=azulprof, draw=azulprof, text=white},
  fl/.style     = {-{Stealth[length=2.6mm]}, draw=vinox, line width=0.8pt},
  flt/.style    = {-{Stealth[length=2.6mm]}, draw=tealx, line width=0.8pt},
  etq/.style    = {font=\scriptsize\itshape, text=azulmedio}
}

\begin{document}

%=======================================================================
\begin{center}
  {\color{tealx}\bfseries\footnotesize SESIÓN 12}\\[4pt]
  {\color{azulprof}\LARGE\bfseries Virtualización y contenedores en Linux}\\[6pt]
  {\color{azulmedio}\small Máquinas virtuales · hipervisores · namespaces · cgroups · redes y persistencia}\\[3pt]
  {\color{azulmedio}\footnotesize Sistemas Operativos Aplicados a las Telecomunicaciones · UNMSM}
\end{center}
\vspace{0.4em}
{\color{grislinea}\hrule height 1pt}
\vspace{1.2em}

%=======================================================================
\section{Punto de partida}

\subsection*{Continuidad del curso: del proceso al aislamiento}

Esta sesión se apoya en lo ya estudiado y no introduce un tema desconectado:

\begin{center}
\begin{tikzpicture}[node distance=5mm]
  \node[cazul, text width=2.9cm] (a) {\textbf{S3--S4}\\[2pt]Sistema de archivos, procesos y CPU};
  \node[cazul, text width=2.9cm, right=of a] (b) {\textbf{S5--S6}\\[2pt]IPC, memoria, buffers y E/S};
  \node[cverde, text width=2.9cm, right=of b] (c) {\textbf{S7}\\[2pt]Sockets, TCP/UDP y servicios};
  \node[crosa, text width=2.9cm, right=of c] (d) {\textbf{S12}\\[2pt]VM, aislamiento y contenedores};
  \draw[fl] (a)--(b); \draw[fl] (b)--(c); \draw[fl] (c)--(d);
\end{tikzpicture}
\end{center}

\begin{clave}
\textbf{Pregunta que organiza la sesión:} ¿cómo puede el sistema operativo presentar
recursos y entornos separados a procesos diferentes?
\end{clave}

\subsection*{Resultados de aprendizaje}

Al finalizar, el estudiante podrá justificar y diagnosticar, no solo repetir definiciones:

\begin{enumerate}[label=\textbf{\arabic*.}]
  \item \textbf{Interpretar} la arquitectura de un hipervisor y de una máquina virtual.
  \item \textbf{Diferenciar} kernel anfitrión, sistema invitado y proceso de contenedor.
  \item \textbf{Explicar} namespaces, cgroups, almacenamiento y redes virtuales.
  \item \textbf{Diagnosticar} evidencia básica de Linux y decidir entre VM y contenedor.
\end{enumerate}

\subsection*{Problema inicial: un servidor debe alojar varios servicios}

Un único servidor físico debe sostener tres servicios con dependencias distintas:
un \textbf{servidor web} (puerto 80), una \textbf{base de datos} (datos persistentes)
y un servicio de \textbf{señalización} (bibliotecas y versiones propias).

\begin{clave}
El problema no es solo ejecutar tres programas: también hay que \textbf{controlar fallas,
recursos y dependencias} sobre un mismo conjunto de hardware.
\end{clave}

\subsection*{Qué controla realmente Linux}

Antes de virtualizar conviene fijar la arquitectura base. Una misma máquina dispone de
hardware físico y de un kernel que \emph{arbitra} el acceso a ese hardware.

\begin{center}
\begin{tikzpicture}[node distance=6mm]
  \node[cazul, text width=11.5cm] (app) {\textbf{Aplicaciones y procesos:} web · monitor · base de datos};
  \node[cverde, text width=11.5cm, below=of app] (k)
    {\textbf{Kernel Linux} \;\textemdash\; planifica CPU · gestiona memoria · controla E/S · aplica permisos};
  \node[caja, below left=12mm and 3.6cm of k.south] (cpu) {CPU};
  \node[caja, right=5mm of cpu] (ram) {RAM};
  \node[caja, right=5mm of ram] (dsk) {Disco};
  \node[caja, right=5mm of dsk] (net) {Red};
  \draw[fl] (app)--(k);
  \draw[fl] (k.south) -- ++(0,-4mm) -| (cpu);
  \draw[fl] (k.south) -- ++(0,-4mm) -| (ram);
  \draw[fl] (k.south) -- ++(0,-4mm) -| (dsk);
  \draw[fl] (k.south) -- ++(0,-4mm) -| (net);
\end{tikzpicture}
\end{center}

\begin{clave}
La virtualización y los contenedores \textbf{no inventan recursos}: reorganizan la manera
en que el kernel —o un sistema invitado— presenta y usa CPU, memoria, disco y red.
\end{clave}

%=======================================================================
\bloque{Bloque 1 · Máquinas virtuales}{Del hardware real a las máquinas virtuales}

\section{Qué significa virtualizar una máquina}

Virtualizar es \textbf{presentar a un sistema invitado una plataforma de hardware virtual}:
CPU, memoria, disco y tarjeta de red que el invitado percibe como propios, aunque sus
recursos provienen del equipo físico.

\begin{center}
\begin{tikzpicture}[node distance=9mm]
  \node[caja, text width=2.6cm] (hw) {\textbf{Hardware físico}\\[2pt]\scriptsize CPU · RAM · NIC};
  \node[cosc, right=16mm of hw, text width=2.4cm] (hv) {\textbf{Hipervisor}};
  \node[cazul, above right=4mm and 16mm of hv, text width=3.2cm] (vm1) {\textbf{VM 1} · Linux};
  \node[cazul, below right=4mm and 16mm of hv, text width=3.2cm] (vm2) {\textbf{VM 2} · otro SO};
  \draw[fl] (hw)--(hv);
  \draw[flt] (hv)--(vm1);
  \draw[flt] (hv)--(vm2);
\end{tikzpicture}
\end{center}

\begin{defbox}
\textbf{SO invitado (\emph{guest}):} sistema operativo que se ejecuta \emph{dentro} de la
máquina virtual, con su propio kernel, sus servicios y sus procesos.\\
\textbf{SO anfitrión (\emph{host}):} sistema que posee el hardware real y sostiene al hipervisor.
\end{defbox}

\begin{clave}
Las CPU, la RAM y la NIC de una VM son \textbf{abstracciones}; sus recursos proceden del
equipo físico y están sujetos a la planificación del anfitrión.
\end{clave}

\subsection{Hipervisores de tipo 1 y tipo 2}

La clasificación describe \textbf{dónde se sitúa la capa de virtualización} respecto del hardware.

\begin{center}
\begin{tikzpicture}[node distance=3mm]
  % Tipo 1
  \node[font=\bfseries\small, text=vinox] (t1) {TIPO 1 · nativo o \emph{bare metal}};
  \node[cazul, below=5mm of t1, text width=5.4cm] (a1) {VM A \quad + \quad VM B};
  \node[crosa, below=of a1, text width=5.4cm] (a2) {Hipervisor / capa de virtualización};
  \node[caja,  below=of a2, text width=5.4cm] (a3) {Hardware físico};
  % Tipo 2
  \node[font=\bfseries\small, text=tealx, right=22mm of t1] (t2) {TIPO 2 · alojado};
  \node[cazul,  below=5mm of t2, text width=5.4cm] (b1) {VM A \quad + \quad VM B};
  \node[crosa,  below=of b1, text width=5.4cm] (b2) {Hipervisor alojado};
  \node[cverde, below=of b2, text width=5.4cm] (b3) {SO anfitrión};
  \node[caja,   below=of b3, text width=5.4cm] (b4) {Hardware físico};
\end{tikzpicture}
\end{center}

\begin{itemize}
  \item \textbf{Tipo 1:} la capa de virtualización se ejecuta directamente sobre el hardware.
        Es el modelo típico de servidores e infraestructura de producción.
  \item \textbf{Tipo 2:} el hipervisor se ejecuta como aplicación sobre un sistema operativo
        anfitrión ya instalado. Es el modelo habitual de escritorio y laboratorio.
\end{itemize}

\begin{aviso}
\textbf{Precisión importante:} KVM en Linux no encaja en una \textquotedblleft caja\textquotedblright{} única.
Requiere explicar la colaboración entre el \textbf{kernel anfitrión} (que aporta las funciones
de virtualización), \textbf{QEMU} (que provee la máquina y sus dispositivos) y la
\textbf{aceleración por hardware} del procesador.
\end{aviso}

\subsection{Una máquina virtual tiene su propio kernel}

El invitado inicia servicios y procesos como si fuese una computadora independiente.
Las capas, de arriba hacia abajo, son:

\begin{center}
\begin{tikzpicture}[node distance=3mm]
  \node[cazul,  text width=11.5cm] (l1) {\textbf{Aplicación invitada} \hfill \scriptsize ejecución en espacio de usuario};
  \node[cverde, below=of l1, text width=11.5cm] (l2) {\textbf{SO invitado (kernel)} \hfill \scriptsize planificación, memoria, sistema de archivos};
  \node[cambar, below=of l2, text width=11.5cm] (l3) {\textbf{CPU / RAM / NIC virtuales} \hfill \scriptsize dispositivos vistos por el invitado};
  \node[crosa,  below=of l3, text width=11.5cm] (l4) {\textbf{QEMU + KVM} \hfill \scriptsize virtualización, aceleración y dispositivos};
  \node[caja,   below=of l4, text width=11.5cm] (l5) {\textbf{Host Linux + hardware} \hfill \scriptsize recursos físicos reales};
  \draw[fl] (l1)--(l2); \draw[fl] (l2)--(l3); \draw[fl] (l3)--(l4); \draw[fl] (l4)--(l5);
\end{tikzpicture}
\end{center}

\begin{clave}
La VM puede \textbf{instalar paquetes y modificar su kernel} sin convertirlo en el kernel
del host: hay una separación real del sistema operativo invitado.
\end{clave}

\subsection{Linux: responsabilidades de KVM, QEMU y libvirt}

Son tres piezas \emph{diferentes} de una misma solución de virtualización, y confundirlas
es uno de los errores más frecuentes.

\begin{center}
\begin{tikzpicture}[node distance=10mm]
  \node[cazul, text width=3.4cm] (lv) {\textbf{libvirt}\\[3pt]\scriptsize Gestión: definir, iniciar, detener y administrar máquinas};
  \node[cverde, right=of lv, text width=3.4cm] (qe) {\textbf{QEMU}\\[3pt]\scriptsize Proceso que proporciona la máquina y los dispositivos};
  \node[cambar, right=of qe, text width=3.4cm] (kv) {\textbf{KVM}\\[3pt]\scriptsize Funciones del kernel que aceleran la CPU virtual};
  \draw[fl] (lv)--(qe); \draw[fl] (qe)--(kv);
\end{tikzpicture}
\end{center}

\begin{defbox}
\cmd{/dev/kvm} es la \textbf{interfaz del kernel} utilizada por el software de virtualización
para solicitar las funciones de aceleración.
\end{defbox}

\begin{errorbox}
\textbf{No confundir:} libvirt \textbf{no ejecuta por sí solo} instrucciones del invitado;
coordina herramientas como QEMU/KVM.
\end{errorbox}

\subsection{Qué ocurre al iniciar una máquina virtual}

No basta con \textquotedblleft abrir un programa\textquotedblright: el hipervisor prepara recursos y luego
arranca al invitado.

\begin{center}
\begin{tikzpicture}[node distance=4.5mm]
  \node[cazul, text width=2.3cm] (p1) {\textbf{1. Configuración}\\[2pt]\scriptsize vCPU, RAM, disco, red};
  \node[cverde, right=of p1, text width=2.3cm] (p2) {\textbf{2. Creación}\\[2pt]\scriptsize proceso QEMU y recursos};
  \node[cazul, right=of p2, text width=2.3cm] (p3) {\textbf{3. Arranque}\\[2pt]\scriptsize firmware y cargador};
  \node[cverde, right=of p3, text width=2.3cm] (p4) {\textbf{4. Kernel invitado}\\[2pt]\scriptsize inicializa servicios};
  \node[cazul, right=of p4, text width=2.3cm] (p5) {\textbf{5. Aplicaciones}\\[2pt]\scriptsize procesos del invitado};
  \draw[fl] (p1)--(p2); \draw[fl] (p2)--(p3); \draw[fl] (p3)--(p4); \draw[fl] (p4)--(p5);
\end{tikzpicture}
\end{center}

\begin{aviso}
El arranque de una VM \textbf{no es instantáneo}: debe iniciar su propio sistema operativo
invitado. Esta es una diferencia operativa decisiva frente a los contenedores.
\end{aviso}

\section{Recursos virtualizados}

\subsection{CPU virtual: ¿cuatro vCPU equivalen a cuatro núcleos?}

\textbf{No necesariamente.} El planificador del host asigna tiempo de CPU física a las vCPU,
igual que lo haría con cualquier otro hilo de ejecución.

\begin{center}
\begin{tikzpicture}[node distance=7mm]
  \node[cazul] (v1) {vCPU 1};
  \node[cazul, right=of v1] (v2) {vCPU 2};
  \node[cazul, right=of v2] (v3) {vCPU 3};
  \node[cazul, right=of v3] (v4) {vCPU 4};
  \node[cverde, below=13mm of $(v2)!0.5!(v3)$, text width=9cm] (sch)
    {\textbf{Planificador del host Linux + KVM}};
  \node[caja, below left=13mm and 2.6cm of sch.south] (c1) {CPU};
  \node[caja, right=10mm of c1] (c2) {CPU};
  \node[caja, right=10mm of c2] (c3) {CPU};
  \draw[fl] (v1)--(sch); \draw[fl] (v2)--(sch); \draw[fl] (v3)--(sch); \draw[fl] (v4)--(sch);
  \draw[flt] (sch.south) -- ++(0,-4mm) -| (c1);
  \draw[flt] (sch.south) -- ++(0,-4mm) -| (c2);
  \draw[flt] (sch.south) -- ++(0,-4mm) -| (c3);
\end{tikzpicture}
\end{center}

\begin{clave}
\textbf{vCPU = unidad de ejecución virtual.} No es una reserva exclusiva de núcleo, salvo
que se configure y se garantice explícitamente (fijación de CPU, reservas, aislamiento).
\end{clave}

\subsection{Memoria virtualizada: asignar no crea RAM}

El anfitrión debe sostener el uso real de sus invitados \emph{además} de su propio consumo.
Ejemplo didáctico con un host de 16 GiB de RAM:

\begin{center}
\begin{tikzpicture}[x=1cm,y=1cm]
  \def\h{0.95}
  \fill[azulprof] (0,0) rectangle (1.8,\h);
  \fill[tealx]    (1.8,0) rectangle (6.6,\h);
  \fill[azulmedio](6.6,0) rectangle (11.4,\h);
  \fill[ambarx]   (11.4,0) rectangle (13.0,\h);
  \node[white,font=\scriptsize\bfseries,align=center] at (0.9,\h/2) {Host\\2 GiB};
  \node[white,font=\small\bfseries,align=center]      at (4.2,\h/2) {VM A · 6 GiB};
  \node[white,font=\small\bfseries,align=center]      at (9.0,\h/2) {VM B · 6 GiB};
  \node[white,font=\scriptsize\bfseries,align=center] at (12.2,\h/2){Libre\\2 GiB};
  \draw[grislinea] (0,0) rectangle (13.0,\h);
\end{tikzpicture}
\end{center}

\begin{itemize}
  \item \textbf{RAM presentada:} configuración nominal del invitado.
  \item \textbf{RAM efectivamente usada:} depende de las aplicaciones y de la actividad real.
  \item \textbf{Presión de memoria:} si la suma de usos reales supera lo disponible,
        puede afectar al host y a las demás VM.
\end{itemize}

\subsection{Disco virtual: un archivo puede representar un disco}

El invitado ve bloques virtuales; el host guarda esa información en almacenamiento físico,
normalmente como un \textbf{archivo de imagen} en formato \cmd{qcow2} o \cmd{raw}.

\begin{center}
\begin{tikzpicture}[node distance=14mm]
  \node[cosc, text width=2.8cm] (d1) {\textbf{Disco de VM}\\[2pt]\scriptsize bloques vistos por el invitado};
  \node[cazul, right=of d1, text width=3.4cm] (d2) {\textbf{Imagen de disco}\\[2pt]\cmd{qcow2} / \cmd{raw}};
  \node[cverde, right=of d2, text width=3.0cm] (d3) {\textbf{Almacenamiento del host}};
  \draw[fl] (d1)--(d2); \draw[fl] (d2)--(d3);
\end{tikzpicture}
\end{center}

\begin{aviso}
Una \textbf{instantánea} (\emph{snapshot}) $\neq$ copia de seguridad completa.
Requiere política de respaldo y \textbf{prueba de recuperación}.
\end{aviso}

\subsection{Red de una VM: NIC virtual, puente y NIC física}

La máquina invitada puede tener una interfaz de red propia sin poseer una tarjeta física.

\begin{center}
\begin{tikzpicture}[node distance=6mm]
  \node[cosc, text width=2.0cm] (vm) {\textbf{VM}\\\scriptsize invitado};
  \node[cverde, right=of vm, text width=1.8cm] (vnic) {\textbf{vNIC}};
  \node[cazul, right=of vnic, text width=2.6cm] (br) {\textbf{Bridge / vSwitch}};
  \node[caja, right=of br, text width=2.0cm] (nic) {\textbf{NIC física}};
  \node[cverde, right=of nic, text width=2.0cm] (ext) {\textbf{Red externa}};
  \draw[fl] (vm)--(vnic); \draw[fl] (vnic)--(br); \draw[fl] (br)--(nic); \draw[fl] (nic)--(ext);
\end{tikzpicture}
\end{center}

En otra configuración, la salida puede usar \textbf{NAT} en lugar de una unión de puente directa.

\begin{defbox}
\textbf{NIC} = interfaz de red física. \textbf{vNIC} = interfaz presentada por el hipervisor.
Crear una vNIC \textbf{no añade} una tarjeta física al servidor.
\end{defbox}

\subsection{Caso de telecomunicaciones: servicios en tres VM}

Separación operativa de roles sobre un mismo servidor físico:

\begin{center}
\begin{tikzpicture}[node distance=8mm]
  \node[cazul,  text width=3.4cm] (g) {\textbf{VM de gestión}\\[2pt]\scriptsize monitoreo y alertas};
  \node[cverde, right=of g, text width=3.4cm] (s) {\textbf{VM de señalización}\\[2pt]\scriptsize servicios de control};
  \node[cambar, right=of s, text width=3.4cm] (d) {\textbf{VM de datos}\\[2pt]\scriptsize persistencia y reportes};
  \node[cosc, below=12mm of s, text width=11.5cm] (base)
    {\textbf{Capa de virtualización + Linux anfitrión + hardware físico}};
  \draw[fl] (g)--(base); \draw[fl] (s)--(base); \draw[fl] (d)--(base);
\end{tikzpicture}
\end{center}

\begin{errorbox}
Separar funciones en VM facilita administrar sus entornos, pero \textbf{no garantiza por sí
solo alta disponibilidad}: si el servidor físico falla, caen las tres máquinas.
\end{errorbox}

\subsection{Qué resuelve y qué no resuelve una VM}

\begin{center}
\begin{tabularx}{\textwidth}{@{}>{\raggedright\arraybackslash}X >{\raggedright\arraybackslash}X@{}}
\toprule
\textbf{\color{tealx}A favor} & \textbf{\color{vinox}Precauciones} \\
\midrule
SO invitado propio, independiente del host & Arranque completo del invitado (tiempo y recursos) \\
Dependencias y bibliotecas separadas & Memoria y servicios duplicados en cada VM \\
Políticas de seguridad y respaldo por VM & Infraestructura física compartida: un fallo afecta a todas \\
\bottomrule
\end{tabularx}
\end{center}

\begin{clave}
Una VM \textbf{no crea recursos}: redistribuye y protege recursos existentes mediante
mecanismos del hipervisor y del host.
\end{clave}

%=======================================================================
\bloque{Bloque 2 · Contenedores}{Aislamiento de procesos dentro de un kernel Linux}

\section{Un contenedor Linux no arranca otro kernel}

Un contenedor ejecuta \textbf{procesos aislados} usando los servicios del \emph{mismo}
kernel anfitrión. Esta es la diferencia esencial con una máquina virtual.

\begin{center}
\begin{tikzpicture}[node distance=3mm]
  \node[font=\bfseries\small, text=azulprof] (h1) {MÁQUINAS VIRTUALES};
  \node[cazul, below=5mm of h1, text width=5.6cm] (m1) {Aplicaciones de VM A / VM B};
  \node[cverde, below=of m1, text width=5.6cm] (m2) {\textbf{Kernel invitado A / Kernel invitado B}};
  \node[crosa, below=of m2, text width=5.6cm] (m3) {Hipervisor + host};
  \node[caja,  below=of m3, text width=5.6cm] (m4) {Hardware};

  \node[font=\bfseries\small, text=azulprof, right=20mm of h1] (h2) {CONTENEDORES LINUX};
  \node[cazul, below=5mm of h2, text width=5.6cm] (n1) {Procesos C1 / Procesos C2};
  \node[crosa, below=of n1, text width=5.6cm] (n2) {Namespaces + cgroups};
  \node[cverde, below=of n2, text width=5.6cm] (n3) {\textbf{ÚNICO kernel Linux del host}};
  \node[caja,  below=of n3, text width=5.6cm] (n4) {Hardware};
\end{tikzpicture}
\end{center}

\begin{clave}
En una VM hay un sistema operativo invitado; en un contenedor Linux habitual
\textbf{no hay kernel invitado}.
\end{clave}

\subsection{Qué significa \textquotedblleft contener\textquotedblright{} un proceso}

Contener significa dos cosas simultáneas: \textbf{aislar su vista} del sistema y
\textbf{controlar los recursos} que puede consumir.

\begin{center}
\begin{tikzpicture}[node distance=13mm]
  \node[cosc, text width=2.4cm] (p) {\textbf{Proceso Linux}};
  \node[cverde, right=of p, text width=3.6cm] (ns) {\textbf{namespaces}\\[2pt]\emph{¿qué puede ver?}};
  \node[cambar, right=of ns, text width=3.6cm] (cg) {\textbf{cgroups}\\[2pt]\emph{¿cuánto puede usar?}};
  \draw[fl] (p)--(ns); \draw[fl] (ns)--(cg);
\end{tikzpicture}
\end{center}

\begin{errorbox}
Juntos permiten construir un entorno de ejecución separado, pero \textbf{NO equivalen a una
VM completa}: el kernel sigue siendo uno solo y es compartido.
\end{errorbox}

\section{Namespaces: distintas vistas sobre el mismo host}

Cada tipo de namespace aísla una \textbf{clase concreta} de recursos.

\begin{center}
\renewcommand{\arraystretch}{1.35}
\begin{tabularx}{\textwidth}{@{}>{\bfseries}l >{\raggedright\arraybackslash}X@{}}
\toprule
\textbf{\normalfont\color{azulprof} Namespace} & \textbf{\color{azulprof} Qué aísla} \\
\midrule
PID   & Identificadores y jerarquía de procesos \\
NET   & Interfaces, rutas, puertos y pila de red \\
MNT   & Puntos de montaje del sistema de archivos \\
UTS   & Nombre de equipo (\emph{hostname}) \\
USER  & Identidades y mapeo UID/GID \\
IPC   & Objetos de comunicación entre procesos \\
\bottomrule
\end{tabularx}
\end{center}

\begin{clave}
Un namespace \textbf{no entrega más RAM ni más CPU}: solo separa la \emph{vista} de un
recurso determinado. La cantidad disponible la gobiernan los cgroups.
\end{clave}

\subsection{Namespace PID: un proceso puede tener IDs distintos según la vista}

Dentro de un contenedor un proceso puede aparecer como \textbf{PID 1} sin ser el PID 1 del host.

\begin{center}
\begin{tikzpicture}[node distance=3mm]
  \node[font=\bfseries\small, text=azulprof] (t1) {VISTA DEL HOST};
  \node[caja, below=5mm of t1, text width=5.4cm, align=left] (h1) {\textbf{PID 1}\hfill init/systemd};
  \node[caja, below=of h1, text width=5.4cm, align=left] (h2) {\textbf{PID 3482}\hfill proceso de aplicación};
  \node[caja, below=of h2, text width=5.4cm, align=left] (h3) {\textbf{PID 3510}\hfill hijo del proceso};

  \node[font=\bfseries\small, text=azulprof, right=22mm of t1] (t2) {VISTA DEL CONTENEDOR};
  \node[cverde, below=5mm of t2, text width=5.4cm, align=left] (c1) {\textbf{PID 1}\hfill aplicación principal};
  \node[cverde, below=of c1, text width=5.4cm, align=left] (c2) {\textbf{PID 2}\hfill proceso hijo};
  \draw[fl] (h2.east) -- node[etq, above]{mismo proceso} (c1.west);
\end{tikzpicture}
\end{center}

\begin{aviso}
Los números son ilustrativos: un PID visto dentro del contenedor \textbf{no identifica el
mismo número} visto desde el host.
\end{aviso}

\subsection{Namespace NET: interfaces y puertos por entorno}

Cada contenedor puede tener su propia tabla de rutas y su propio espacio de puertos.

\begin{center}
\begin{tikzpicture}[node distance=12mm]
  \node[cazul, text width=3.6cm] (a) {\textbf{CONTENEDOR A}\\[2pt]\cmd{eth0} · 172.18.0.2\\\scriptsize puerto interno 80};
  \node[caja, right=of a, text width=3.0cm] (br) {\textbf{bridge de red virtual}};
  \node[cverde, right=of br, text width=3.6cm] (b) {\textbf{CONTENEDOR B}\\[2pt]\cmd{eth0} · 172.18.0.3\\\scriptsize puerto interno 80};
  \draw[fl] (a)--(br); \draw[fl] (br)--(b);
\end{tikzpicture}
\end{center}

Las IP y los puertos son ilustrativos; las interfaces se suelen enlazar mediante
\textbf{pares veth} conectados a un puente.

\subsection{Namespace MNT: distintas vistas de archivos}

Dos procesos pueden recorrer la ruta \cmd{/app} y encontrar contenidos diferentes.

\begin{center}
\renewcommand{\arraystretch}{1.3}
\begin{tabularx}{\textwidth}{@{}>{\bfseries}l X X@{}}
\toprule
\textbf{\normalfont\color{azulprof} Entorno} & \textbf{\color{azulprof} Ruta observada} & \textbf{\color{azulprof} Contenido real} \\
\midrule
Host          & \cmd{/srv/app-uno}, \cmd{/srv/app-dos} & directorios físicos del anfitrión \\
Contenedor A  & \cmd{/app}, \cmd{/config} & datos A, configuración A \\
Contenedor B  & \cmd{/app}, \cmd{/config} & datos B, configuración B \\
\bottomrule
\end{tabularx}
\end{center}

\begin{errorbox}
Separar puntos de montaje \textbf{no cifra automáticamente} los datos ni sustituye las
autorizaciones del sistema (permisos, propietarios, políticas).
\end{errorbox}

\section{Cgroups: quién controla CPU y memoria}

Un \textbf{cgroup} (\emph{control group}) agrupa procesos y permite \textbf{medir} o
\textbf{restringir} su consumo de recursos. Los archivos de control más usados en
\emph{cgroup v2} son \cmd{cpu.max} para CPU y \cmd{memory.max} para memoria.

\begin{center}
\begin{tikzpicture}[x=1cm,y=1cm]
  \node[font=\small\bfseries,anchor=east] at (-0.2,0.3) {Servicio A};
  \draw[grislinea] (0,0) rectangle (8,0.6);
  \fill[tealx] (0,0) rectangle (5.6,0.6);
  \node[white,font=\small\bfseries] at (2.8,0.3) {70\%};
  \node[font=\small\bfseries,anchor=east] at (-0.2,-0.9) {Servicio B};
  \draw[grislinea] (0,-1.2) rectangle (8,-0.6);
  \fill[azulmedio] (0,-1.2) rectangle (3.2,-0.6);
  \node[white,font=\small\bfseries] at (1.6,-0.9) {40\%};
\end{tikzpicture}
\end{center}

\subsection{Una cuota de CPU no representa una reserva exclusiva}

El control \textbf{limita cuánto puede consumir} un grupo; no garantiza tiempo dedicado.
Supongamos la configuración:

\begin{codigo}
cpu.max = 50000 100000      # cuota (us)  periodo (us)
\end{codigo}

\begin{center}
\begin{tikzpicture}[x=1cm,y=1cm]
  \foreach \i in {0,...,4}{ \fill[tealx] (\i*1.3,0) rectangle (\i*1.3+1.2,0.8); }
  \foreach \i in {5,...,9}{ \fill[grisfondo] (\i*1.3,0) rectangle (\i*1.3+1.2,0.8);
                            \draw[grislinea] (\i*1.3,0) rectangle (\i*1.3+1.2,0.8); }
  \foreach \i/\lab in {0/{0--10},1/{10--20},2/{20--30},3/{30--40},4/{40--50},
                       5/{50--60},6/{60--70},7/{70--80},8/{80--90},9/{90--100}}
    { \node[font=\tiny,text=azulmedio] at (\i*1.3+0.6,-0.25) {\lab}; }
  \node[font=\small\bfseries,text=azulprof,anchor=west] at (0,-0.85) {Periodo de 100\,000 µs};
  \node[font=\small\bfseries,text=azulprof,anchor=east] at (12.7,-0.85) {Cuota de 50\,000 µs};
\end{tikzpicture}
\end{center}

\begin{aviso}
\textbf{Interpretación:} se permite hasta \textbf{50 ms de CPU cada 100 ms} de periodo,
equivalente a \textbf{0,5 CPU} en la cuota. Si el proceso agota la cuota antes de terminar
el periodo, queda \emph{throttled} (limitado) hasta el siguiente periodo.
\end{aviso}

\subsection{La memoria puede limitarse y fallar por exceso}

Si un cgroup no logra reducir su uso al acercarse al límite de \cmd{memory.max},
pueden producirse eventos \textbf{OOM}.

\begin{center}
\begin{tikzpicture}[x=1cm,y=1cm]
  \draw[grislinea] (0,0) rectangle (11,0.9);
  \fill[tealx] (0,0) rectangle (7.6,0.9);
  \draw[vinox,line width=1pt] (7.6,-0.35) -- (7.6,1.25);
  \node[font=\small\bfseries,text=vinox,anchor=south] at (8.9,1.3) {Límite duro · memory.max};
  \node[font=\small\bfseries,text=azulprof,anchor=west] at (0,-0.8)
    {Reclamar memoria \;$\rightarrow$\; si no alcanza: posible OOM del cgroup};
\end{tikzpicture}
\end{center}

\begin{defbox}
\textbf{OOM (\emph{Out Of Memory})} = condición de falta de memoria que puede terminar
procesos del grupo. \textbf{No es idéntico} a consumir el 100\,\% de CPU: son fallas
de naturaleza distinta y se diagnostican con contadores distintos.
\end{defbox}

\subsection{Aislamiento no equivale a seguridad absoluta}

Un contenedor necesita \textbf{defensa en profundidad} además de namespaces y cgroups:

\begin{center}
\renewcommand{\arraystretch}{1.3}
\begin{tabularx}{\textwidth}{@{}>{\bfseries}l >{\raggedright\arraybackslash}X@{}}
\toprule
\textbf{\normalfont\color{azulprof} Capa} & \textbf{\color{azulprof} Aporte} \\
\midrule
Namespaces      & Separación de vistas \\
Cgroups         & Límites y contabilidad de recursos \\
Identidad       & UID no privilegiado y capacidades mínimas \\
Seccomp / MAC   & Restricción de llamadas al sistema y operaciones \\
Actualizaciones & Host, imágenes y aplicaciones al día \\
\bottomrule
\end{tabularx}
\end{center}

\section{Almacenamiento de contenedores}

\subsection{Imagen, contenedor y volumen son objetos diferentes}

\begin{center}
\begin{tikzpicture}[node distance=12mm]
  \node[cazul, text width=3.4cm] (im) {\textbf{IMAGEN}\\[3pt]\scriptsize capas base\\(solo lectura)\\[2pt]\emph{plantilla}};
  \node[cverde, right=of im, text width=3.4cm] (co) {\textbf{CONTENEDOR}\\[3pt]\scriptsize capa escribible\\proceso en ejecución\\[2pt]\emph{instancia}};
  \node[cambar, right=of co, text width=3.4cm] (vo) {\textbf{VOLUMEN}\\[3pt]\scriptsize datos externos a la\\capa escribible\\[2pt]\emph{persistencia}};
  \draw[fl] (im)--(co); \draw[fl] (co)--(vo);
\end{tikzpicture}
\end{center}

\begin{clave}
Al eliminar un contenedor, su \textbf{capa escribible se elimina}; un volumen separado
puede conservar los datos.
\end{clave}

\subsection{Persistencia: la vida de los datos no debe depender del proceso}

Una base de datos necesita almacenamiento \textbf{fuera} de la capa escribible efímera.

\begin{center}
\begin{tikzpicture}[node distance=13mm]
  \node[cosc, text width=3.0cm] (db) {\textbf{Proceso de base de datos}};
  \node[cverde, right=of db, text width=2.8cm] (ct) {\textbf{Contenedor}};
  \node[cambar, right=of ct, text width=3.0cm] (vl) {\textbf{Volumen persistente}};
  \draw[fl] (db)--(ct); \draw[fl] (ct)--(vl);
\end{tikzpicture}
\end{center}

\begin{aviso}
El volumen puede persistir tras eliminar el contenedor, \textbf{pero sigue necesitando
copias de seguridad}: persistencia no es respaldo.
\end{aviso}

\section{Docker: de un comando a un proceso aislado}

Docker es una \textbf{plataforma de gestión} que utiliza funciones del kernel Linux;
no es un mecanismo de aislamiento por sí mismo.

\begin{center}
\begin{tikzpicture}[node distance=7mm]
  \node[cazul, text width=2.6cm] (cl) {\textbf{Cliente}\\[2pt]\cmd{docker run}};
  \node[crosa, right=of cl, text width=2.6cm] (ap) {\textbf{API / daemon}\\[2pt]\cmd{dockerd}};
  \node[cverde, right=of ap, text width=2.6cm] (rt) {\textbf{Runtime}\\[2pt]\scriptsize crea la instancia};
  \node[cambar, right=of rt, text width=2.9cm] (ke) {\textbf{Kernel Linux}\\[2pt]\scriptsize namespaces + cgroups};
  \draw[fl] (cl)--(ap); \draw[fl] (ap)--(rt); \draw[fl] (rt)--(ke);
\end{tikzpicture}
\end{center}

\begin{clave}
La terminal \textbf{no \textquotedblleft es\textquotedblright{} el contenedor}: envía una solicitud a componentes que
crean y gestionan procesos.
\end{clave}

\subsection{Red Docker: ¿qué significa \texttt{-p 8080:80}?}

Publica el puerto 80 del contenedor a través del puerto 8080 del host.

\begin{center}
\begin{tikzpicture}[node distance=14mm]
  \node[caja, text width=2.6cm] (cli) {\textbf{Cliente de red}};
  \node[cazul, right=of cli, text width=3.4cm] (host) {\textbf{HOST LINUX}\\[2pt]\cmd{IP\_HOST:8080}};
  \node[cverde, right=of host, text width=3.4cm] (cont) {\textbf{CONTENEDOR}\\[2pt]\cmd{IP\_CONT:80}};
  \draw[fl] (cli)--(host); \draw[fl] (host)--(cont);
\end{tikzpicture}
\end{center}

\begin{defbox}
\cmd{-p 8080:80} $\rightarrow$ \textbf{puerto publicado del host : puerto del contenedor}.
\end{defbox}

\begin{errorbox}
Publicar un puerto es un \textbf{mecanismo de acceso de red}; \textbf{no modifica} el número
de puerto en el que escucha el proceso interno.
\end{errorbox}

\subsection{Interpretar \texttt{docker run}}

\begin{codigo}
docker run --rm -d --name web -p 8080:80 nginx:alpine
\end{codigo}

\begin{center}
\renewcommand{\arraystretch}{1.3}
\begin{tabularx}{\textwidth}{@{}>{\ttfamily}l >{\raggedright\arraybackslash}X@{}}
\toprule
\textbf{\normalfont\color{azulprof} Elemento} & \textbf{\color{azulprof} Significado} \\
\midrule
nginx:alpine & imagen de ejemplo (plantilla de la que se crea la instancia) \\
-{}-name web & nombre de instancia, para referirse al contenedor \\
-d           & ejecución desacoplada (en segundo plano) \\
-{}-rm       & eliminar el contenedor al finalizar \\
-p 8080:80   & publicación de puerto host:contenedor \\
\bottomrule
\end{tabularx}
\end{center}

\begin{aviso}
Requiere descarga de la imagen si no existe en caché; ejecutar \textbf{solo en un entorno
autorizado}.
\end{aviso}

\subsection{Comandos Docker de inspección}

\begin{center}
\renewcommand{\arraystretch}{1.3}
\begin{tabularx}{\textwidth}{@{}>{\ttfamily}l >{\raggedright\arraybackslash}X@{}}
\toprule
\textbf{\normalfont\color{azulprof} Comando} & \textbf{\color{azulprof} Qué informa} \\
\midrule
docker version      & identifica cliente y servidor disponibles \\
docker images       & lista las imágenes almacenadas localmente \\
docker ps           & lista los contenedores en ejecución \\
docker inspect web  & presenta configuración y estado de \textquotedblleft web\textquotedblright \\
docker logs web     & muestra la salida registrada de la aplicación \\
\bottomrule
\end{tabularx}
\end{center}

Si Docker no está instalado, \textbf{interpretar los comandos y sus salidas simuladas es
suficiente} para los fines de la sesión.

%=======================================================================
\bloque{Bloque 3 · Aplicación y diagnóstico}{Interpretación de evidencia Linux y decisiones de arquitectura}

\section{Observar el sistema sin modificarlo}

\subsection{Lectura de namespaces y cgroups}

Estos comandos \textbf{inspeccionan} el sistema; sus resultados dependen del equipo:

\begin{codigo}
$ uname -r
$ lsns -t pid,net
$ readlink /proc/self/ns/net
$ cat /proc/self/cgroup
\end{codigo}

\begin{center}
\renewcommand{\arraystretch}{1.3}
\begin{tabularx}{\textwidth}{@{}>{\ttfamily\bfseries}l >{\raggedright\arraybackslash}X@{}}
\toprule
\textbf{\normalfont\color{azulprof} Comando} & \textbf{\color{azulprof} Interpretación} \\
\midrule
uname -r  & versión del kernel que ejecuta Linux \\
lsns      & namespaces visibles, según permisos del usuario \\
readlink  & identificador del namespace del proceso actual \\
/proc     & información expuesta por el kernel sobre procesos y cgroups \\
\bottomrule
\end{tabularx}
\end{center}

\begin{errorbox}
\textbf{No copie salidas inventadas como evidencia real:} interprete las que devuelva cada
instalación concreta.
\end{errorbox}

\subsection{Diagnóstico: ¿estamos dentro de una VM?}

Tres señales orientan el análisis; \textbf{ninguna aislada prueba todo el entorno}.

\begin{center}
\renewcommand{\arraystretch}{1.3}
\begin{tabularx}{\textwidth}{@{}>{\ttfamily\bfseries}l >{\raggedright\arraybackslash}X@{}}
\toprule
\textbf{\normalfont\color{azulprof} Señal} & \textbf{\color{azulprof} Qué aporta} \\
\midrule
systemd-detect-virt & puede identificar ciertas tecnologías de virtualización \\
lscpu               & enumera propiedades de CPU; los campos dependen del hardware y del sistema \\
/dev/kvm            & su existencia y permisos ayudan a examinar el soporte de KVM \\
\bottomrule
\end{tabularx}
\end{center}

\begin{aviso}
\textbf{No confundir} \emph{soporte} de virtualización con una VM ya creada o en ejecución:
la presencia de \cmd{/dev/kvm} indica capacidad del host, no que usted esté dentro de un invitado.
\end{aviso}

\section{Decidir entre VM y contenedor}

\subsection{Caso A: aplicación heredada que exige otro kernel}

\textbf{Requisito principal:} ejecutar software que depende de un sistema operativo invitado
distinto al kernel del host.

\begin{defbox}
\textbf{Opción razonable: máquina virtual} con SO invitado compatible.\\
\emph{Justificación:} permite un entorno invitado completo, con su kernel, sus servicios y
sus dependencias propias; un contenedor Linux no puede aportar otro kernel.
\end{defbox}

\subsection{Caso B: varios microservicios Linux independientes}

Tres componentes —\textbf{API}, \textbf{procesador} y \textbf{notificaciones}— usan software
compatible con el kernel anfitrión y requieren despliegue frecuente.

\begin{center}
\begin{tikzpicture}[node distance=14mm]
  \node[cverde, text width=2.6cm] (api) {\textbf{API}};
  \node[cverde, right=of api, text width=2.6cm] (pro) {\textbf{Procesador}};
  \node[cverde, right=of pro, text width=2.6cm] (not) {\textbf{Notificaciones}};
  \node[cazul, below=11mm of pro, text width=11.5cm] (ker)
    {\textbf{Kernel Linux compartido + aislamiento y límites por servicio}};
  \draw[fl] (api)--(ker); \draw[fl] (pro)--(ker); \draw[fl] (not)--(ker);
\end{tikzpicture}
\end{center}

\begin{defbox}
\textbf{Candidato: contenedores}, con control de recursos (cgroups), seguridad, red y
persistencia definidas según el servicio.
\end{defbox}

\subsection{Comparación técnica: VM y contenedor}

La elección depende de \textbf{requisitos}, no de una etiqueta de \textquotedblleft mejor tecnología\textquotedblright.

\begin{center}
\renewcommand{\arraystretch}{1.35}
\begin{tabularx}{\textwidth}{@{}>{\bfseries}l >{\raggedright\arraybackslash}X >{\raggedright\arraybackslash}X@{}}
\toprule
\textbf{\normalfont\color{azulprof} Dimensión} & \textbf{\color{azulprof} Máquina virtual} & \textbf{\color{azulprof} Contenedor Linux} \\
\midrule
Kernel       & Propio del SO invitado          & Compartido con el host Linux \\
Inicio       & Arranque de SO invitado         & Inicio de procesos aislados \\
Recursos     & Dispositivos / CPU / RAM virtuales & Recursos del host regulados \\
Dependencias & Puede incorporar otro SO        & Espacio de usuario compatible \\
Aislamiento  & Frontera de virtualización      & Frontera basada en kernel y políticas \\
Datos        & Disco / volumen de VM           & Capa escribible y volúmenes \\
\bottomrule
\end{tabularx}
\end{center}

\section{Ejercicios resueltos}

\subsection*{Ejercicio 1 · Recursos compartidos}

\textbf{Enunciado.} Un servidor tiene \textbf{8 núcleos físicos y 16 GiB de RAM}.
Dos máquinas virtuales solicitan \textbf{6 vCPU y 6 GiB} cada una.

\begin{center}
\begin{tikzpicture}[node distance=12mm]
  \node[cazul, text width=3.6cm] (h) {\textbf{HOST}\\[3pt]8 núcleos físicos\\16 GiB RAM};
  \node[cverde, right=of h, text width=3.0cm] (a) {\textbf{VM A}\\[3pt]6 vCPU\\6 GiB RAM};
  \node[cverde, right=of a, text width=3.0cm] (b) {\textbf{VM B}\\[3pt]6 vCPU\\6 GiB RAM};
\end{tikzpicture}
\end{center}

\textbf{Resolución.}
\begin{itemize}
  \item \textbf{CPU:} $6 + 6 = 12$ vCPU configuradas frente a 8 núcleos físicos.
        Hay \emph{sobresuscripción}: es admisible porque las vCPU no consumen tiempo
        permanentemente, pero bajo carga simultánea habrá competencia y latencia.
  \item \textbf{Memoria:} $6 + 6 = 12$ GiB asignados a las VM dejan \textbf{4 GiB} para el
        host y su margen operativo. La memoria, a diferencia de la CPU, no se
        \textquotedblleft turna\textquotedblright: si el uso real se acerca al total, aparece presión de memoria.
\end{itemize}

\begin{clave}
\textbf{Conclusión:} la configuración es viable solo si el uso real de CPU es moderado y
si 4 GiB bastan para el anfitrión. Asignar no equivale a disponer.
\end{clave}

\subsection*{Ejercicio 2 · Dos contenedores escuchan en el puerto 80}

\textbf{Enunciado.} El contenedor A se publica con \cmd{-p 8080:80} y el contenedor B con
\cmd{-p 8081:80}. ¿Por qué es válido que ambos usen internamente el puerto 80?

\begin{center}
\begin{tikzpicture}[node distance=5mm]
  \node[caja, text width=3.2cm] (h1) {\textbf{HOST :8080}};
  \node[cazul, right=18mm of h1, text width=4.0cm] (a) {\textbf{Contenedor A :80}};
  \node[caja, below=10mm of h1, text width=3.2cm] (h2) {\textbf{HOST :8081}};
  \node[cverde, right=18mm of h2, text width=4.0cm] (b) {\textbf{Contenedor B :80}};
  \draw[fl] (h1)--(a); \draw[fl] (h2)--(b);
\end{tikzpicture}
\end{center}

\textbf{Resolución.} Los puertos internos pertenecen a \textbf{namespaces de red separados}:
cada contenedor tiene su propia pila y su propio espacio de puertos, de modo que el
\textquotedblleft puerto 80\textquotedblright{} de A y el de B son objetos distintos. En el host, en cambio, el espacio
de puertos es único, y por eso las publicaciones deben usar números diferentes (8080 y 8081).

\subsection*{Ejercicio 3 · Identificar la causa probable}

\textbf{Enunciado.} Una aplicación se ralentiza. Se dispone de la siguiente evidencia:

\begin{codigo}
cpu.stat
  nr_throttled    135
  throttled_usec  950000

memory.events
  oom       0
  oom_kill  0
\end{codigo}

\textbf{Hipótesis inicial.} Hay \textbf{señales de limitación por cuota de CPU}
(\cmd{nr\_throttled} y \cmd{throttled\_usec} crecen), mientras que los contadores de
memoria mostrados \textbf{no indican terminaciones por OOM}. La ralentización es
coherente con \emph{throttling} del cgroup, no con falta de memoria.

\begin{errorbox}
\textbf{No afirmar causalidad solo con un contador:} hay que comparar intervalos de tiempo,
carga, latencia y configuración antes de concluir.
\end{errorbox}

\subsection*{Ejercicio autónomo · Justificar una arquitectura}

Una plataforma de monitoreo debe integrar un componente heredado y cuatro servicios Linux.
Responda y \emph{justifique}:

\begin{enumerate}[label=\textbf{\Alph*.}]
  \item ¿Qué parte conviene evaluar como VM si exige un SO invitado diferente?
  \item ¿Qué componentes podrían ser contenedores si comparten el kernel Linux?
  \item ¿Cómo separaría redes, procesos y límites de recursos?
  \item ¿Dónde almacenaría datos que deben sobrevivir a recreaciones?
  \item ¿Qué comandos de solo lectura usaría para obtener evidencia?
\end{enumerate}

%=======================================================================
\section{Cierre}

\subsection*{Mapa conceptual: ¿dónde se produce el aislamiento?}

Tres preguntas evitan mezclar tecnologías:

\begin{center}
\begin{tikzpicture}[node distance=8mm]
  \node[cazul, text width=3.7cm] (q1) {\textbf{¿Quién tiene el kernel?}\\[3pt]\scriptsize VM: invitado propio\\Contenedor: host común};
  \node[cverde, right=of q1, text width=3.7cm] (q2) {\textbf{¿Quién ve los recursos?}\\[3pt]\scriptsize Namespaces alteran la vista de determinados recursos};
  \node[cambar, right=of q2, text width=3.7cm] (q3) {\textbf{¿Quién controla el consumo?}\\[3pt]\scriptsize Cgroups contabilizan y limitan CPU, memoria y otros};
\end{tikzpicture}
\end{center}

\begin{clave}
Recurso físico compartido $\neq$ vista compartida $\neq$ kernel compartido:
son \textbf{propiedades distintas} y conviene nombrarlas por separado.
\end{clave}

\subsection*{Conclusión de la sesión}

\begin{tcolorbox}[enhanced, colback=azulprof, colframe=azulprof, arc=3pt,
  left=14pt,right=14pt,top=12pt,bottom=12pt]
{\color{white}\large\bfseries Virtualizar es crear un entorno de máquina.\\
Contenerizar es aislar procesos.}\\[6pt]
{\color{white!80!azulmedio}\small Para escoger: kernel requerido, recursos, red, datos,
seguridad y operación.}
\end{tcolorbox}

%=======================================================================
\section*{Fuentes técnicas y lecturas recomendadas}
\addcontentsline{toc}{section}{Fuentes}

\begin{itemize}[leftmargin=1.2em]
  \item \textbf{Linux man-pages — namespaces(7).}\\ \url{https://man7.org/linux/man-pages/man7/namespaces.7.html}
  \item \textbf{Linux Kernel — Control Group v2.}\\ \url{https://www.kernel.org/doc/html/latest/admin-guide/cgroup-v2.html}
  \item \textbf{libvirt — QEMU/KVM hypervisor driver.}\\ \url{https://www.libvirt.org/drvqemu.html}
  \item \textbf{Docker Docs — Docker overview.}\\ \url{https://docs.docker.com/get-started/docker-overview/}
  \item \textbf{Docker Docs — Networking.}\\ \url{https://docs.docker.com/engine/network/}
  \item \textbf{Docker Docs — Volumes.}\\ \url{https://docs.docker.com/engine/storage/volumes/}
\end{itemize}

\end{document}
